\documentclass[10pt,a4paper]{amsart}
\usepackage[margin=19mm]{geometry}
\usepackage{amsmath,amssymb,mathtools}
\usepackage[scr=rsfs,cal=euler]{mathalfa}
\usepackage{xfrac}
\usepackage[unicode,hidelinks,bookmarks=false]{hyperref}
\newcommand{\ie}{i.e.\ }
\newcommand{\cf}{cf.\ }
\newcommand{\resp}{respectively\ }
\newcommand{\Subsection}{\subsection}
\providecommand{\nobreakdash}{\nobreak\hbox{-}}
\setlength{\emergencystretch}{2em}
\multlinegap0pt
\usepackage{mathtools}
\usepackage{amssymb}
\usepackage{enumerate}
\usepackage{tikz}
\newtheorem{lemma}{Lemma}
\def\N{\mathbb{N}}
\def\cL{{\mathcal L}}
\newcommand{\ra}{\rightarrow}
\title{Coloring powers of random graphs}
\author{Alan Frieze, Ross Kang, Aditya Raut, Michelle Sweering, Hilde Verbeek}
\date{Source: arXiv:2604.14006v1 (CC BY 4.0)}
\begin{document}
\maketitle

\noindent\emph{Source and licence.} Coloring powers of random graphs. Version arXiv:2604.14006v1 (CC BY 4.0), distributed by arXiv under the Creative Commons Attribution 4.0 International licence, http://creativecommons.org/licenses/by/4.0/, as recorded in the arXiv licence metadata for this exact version and shown on the abstract page https://arxiv.org/abs/2604.14006v1. Full licence notice: \texttt{licenses/arxiv-2604.14006-CC-BY-4.0.txt}.\par\smallskip

\noindent\textit{Editorial scope.} Counted unit: the article's lemma lem2 (source0.tex lines 352--354). The article's complete original proof of it is delivered with it (source0.tex lines 355--390, one \texttt{proof} environment of 36 lines). No mathematics is written, repaired or re-derived: the statement and the proof are the article's own lines, in the article's own notation, and no retained line is substituted or re-flowed.  No further source-local context is needed: every cross-reference of every retained line resolves inside the two delivered spans.  Nothing was omitted from any retained span, so the package carries no omission record.  No retained line contains a citation, so the wrapper carries no bibliography and no bibliography entry.  No retained line contains a figure environment, an image-inclusion command or drawing code, so no image is delivered and the package declares no asset dependency.  Editorial formatting only, in addition: an AMS \texttt{amsart} preamble that restates the article's own declarations and macros used by the retained lines, namely the environments \texttt{lemma} on separate counters, together with the standard packages those lines need.  The retained environments print in this document as Lemma 1 in order of appearance, whereas the article prints the same items with the numbering of its own full document.  The complete native source of the article as distributed in the arXiv e-print for this exact version is bundled unchanged as \texttt{sources/198511/source0.tex}.
\begin{lemma}\label{lem2}
	For $\ell_0=1$ and $\ell_1,...,\ell_r \in \N$ such that $\sum\limits_{i=1}^r \ell_i = D$, we have $\min \sum\limits_{i=1}^r \ell_i \log \frac{\ell_i}{\ell_{i-1}} \geq D\log_{(r)}D - O(D)$, for sufficiently large $D$.
\end{lemma}

\begin{proof}
	We proceed by induction on $r$. For $r=1$, the result holds since we have $\ell_1=D$, implying that $\ell_1\log\frac{\ell_1}{\ell_0} = D\log D$. Assume that the result holds for $r-1$.  

{\bf Case 1: $(D-\ell_r)\log_{(r-1)}(D-\ell_r) \geq D\log_{(r)} d$:}\\
Because $\sum\limits_{i=1}^{r-1} \ell_i = D-\ell_r$, from the induction hypothesis we have $\sum\limits_{i=1}^{r-1} \ell_i\log\frac{\ell_i}{\ell_{i-1}} \geq (D-\ell_r)\log_{(r-1)}(D-\ell_r)$. So this case is done.

{\bf Case 2: $(D-\ell_r)\log_{(r-1)}(D-\ell_r)< D\log_{(r)} D$:}\\
For $D$ sufficiently large, we have $\frac{D}{10}\log_{(r-1)} \frac{D}{10} > D\log_{(r)} D$. Hence $\ell_r \geq \frac{9D}{10}$, implying $\ell_{r-1} \leq \frac{D}{10}$ and $\frac{\ell_r}{\ell_{r-1}} \geq 9$.

We use Lagrange multipliers. But first we deal with the constraints $\ell_i\geq 0$ for $i=1,2,\ldots,r$. If $\ell_i=0$ and $\ell_{i+1}\neq 0$ then $\sum\limits_{i=1}^r \ell_i \log \frac{\ell_i}{\ell_{i-1}}=\infty$. If $\ell_i=\ell_{i+1}=\cdots=\ell_r=0$ then $\sum\limits_{i=1}^r \ell_i \log \frac{\ell_i}{\ell_{i-1}}=\sum\limits_{i=1}^{i-1} \ell_i \log \frac{\ell_i}{\ell_{i-1}}$ and the result follows by induction.  So in effect there is one constraint: $\sum\limits_{i=1}^r \ell_i = D$.

Define $\cL(\ell_1,...,\ell_r,\lambda) = \sum\limits_{i=1}^r \ell_i \log \frac{\ell_i}{\ell_{i-1}} + \lambda \left(\sum\limits_{i=1}^r \ell_i - D\right)$. By the Lagrange multiplier theorem, the minima must satisfy $\frac{\partial\cL}{\partial \ell_i} = 0$ for all $i$. Notice that $\frac{\partial \cL}{\partial \ell_i} = 1 + \log\frac{\ell_{i}}{\ell_{i-1}} - \frac{\ell_{i+1}}{\ell_i} + \lambda$ for $1 \leq i < r$, and $\frac{\partial\cL}{\partial\ell_r} = 1+\log \frac{\ell_r}{\ell_{r-1}} +\lambda$. Let $p_i := \frac{\ell_i}{\ell_{i-1}}$. From $\frac{\partial\cL}{\partial \ell_r}=0$, we have $1+\log p_r + \lambda = 0$ which implies that $p_r = e^{-(\lambda+1)}$. For $1\leq i<r$, from $\frac{\partial\cL}{\partial\ell_i}=0$ we have $1+\log p_i -p_{i+1}+\lambda = 0$ which  implies that $p_i = e^{p_{i+1}-(1+\lambda)} = p_r\cdot e^{p_{i+1}}\geq p_rp_{i+1}$. Thus, we can iteratively obtain the exact expressions for $p_1,p_2,...,p_{r-1}$ in terms of $p_r$. Now recall that $p_r\geq 9$ from induction hypothesis, hence $p_i \geq 9$ for all $i$. 
	
Since $\ell_0=1$, $\ell_1=p_1$. Now $\ell_i = p_i\cdot p_{i-1}\cdots p_1$, for all $i=1,2,...,r$ and then $\frac{\ell_r}{\ell_i} = p_rp_{r-1}...p_{i+1} \geq 9^{r-i}$. Now $\sum\limits_{i=1}^r \ell_i = D$ and so clearly $D > \ell_r$. Moreover, 
\[
D = \ell_r \left(\sum\limits_{i=1}^r \frac{\ell_i}{\ell_r}\right) \leq \ell_r \left(\sum\limits_{i=1}^r \frac{1}{9^{r-i}}\right) < \ell_r\left(\sum\limits_{j=0}^\infty \frac{1}{9^j}\right) = \frac{9}{8}\cdot \ell_r. 
\]
We can thus assume that $D = c_1 \cdot \ell_r$ for some $c_1 \in (1,9/8)$. So
\begin{align*}
D= c_1 \cdot p_r \cdot p_{r-1} \cdots p_1 & = c_1 \cdot p_r \cdot \left(p_r e^{p_r}\right) \cdot\left(p_r e^{p_r e^{p_r}}\right) \cdots \left(p_r \underbrace{e^{p_r e^{... e^{p_r}}}}_{\substack{\text{exponential tower}\\\text{of height }r-1}}\right)\\
		& = c_1 \cdot p_r^r \cdot \left(e^{p_r} e^{p_r e^{p_r}} \cdots \underbrace{e^{p_r e^{... e^{p_r}}}}_{\substack{\text{exponential tower}\\\text{of height }r-1}} \right).\\
\end{align*}
Applying the logarithmic function to both sides $(r-1)$ times, we have $p_r\leq \log_{(r-1)} D = p_r + O(\log p_r)$, implying that $\log_{(r-1)}D-c\log_{(r)}D\leq p_r\leq \log_{(r-1)}D$ for some constant $c>0$. We see from the above that $p_{r-1}=p_re^{p_r}$ and that $p_i\geq p_rp_{i+1}$. It follows that $\ell_i\leq \ell_rp_r^{i-r}$ and so 
\[
\ell_r = D(1-\eta)\text{ for }\eta \leq \frac{2}{\log_{(r-1)}D} \ra 0.
\]
 Let $T_i : = \ell_i\log\frac{\ell_i}{\ell_{i-1}}$. Then we have 
\[
\frac{T_i}{T_{i-1}}=p_i\frac{\log p_i}{\log p_{i-1}}=\frac{p_i\log p_i}{p_i+\log p_r}\geq \tfrac12\log p_i>1.
\]
It follows that for all $i<r-1$, we have $T_i <T_{r-1}$. Now $\log p_{r-1} = p_r + \log p_r$ implies that 
\[
T_{r-1} = \ell_{r-1}\log p_{r-1}=\ell_{r-1}(p_r + \log p_r)\leq \frac{\ell_r}{p_r}(p_r + \log p_r)  = O(\ell_r) = O(D).
\]
Thus, the objective is dominated by last summand $T_r = \ell_r\log\frac{\ell_r}{\ell_{r-1}}$, resulting in minimum value of at least $D\log_{(r)}D+O(D)$.
 \end{proof}

\end{document}
